\documentclass[10pt,a4paper]{article}
\usepackage[utf8]{inputenc}
\usepackage[T1]{fontenc}
\usepackage[english]{babel}
\usepackage{geometry}
\geometry{margin=1.8cm}
\usepackage{enumitem}
\usepackage{booktabs}
\usepackage{fancyhdr}
\usepackage{lastpage}
\usepackage{tcolorbox}
\usepackage{hyperref}

\pagestyle{fancy}
\fancyhf{}
\rfoot{\thepage/\pageref{LastPage}}
\lhead{\small TOEFL ITP 2026 --- Day 9}
\rhead{\small MPrometeo}
\renewcommand{\headrulewidth}{0.4pt}
\setlength{\parindent}{0pt}
\setlength{\parskip}{5pt}

\begin{document}

\begin{center}
{\LARGE \textbf{DAY 9 --- Saturday Mar 7}}\\[4pt]
{\large Grammar Intensive: Relative Clauses + Inference}\\[2pt]
{\normalsize Theme: The Pemex Crisis --- From C\'ardenas to Huachicol}\\[2pt]
\rule{\textwidth}{0.8pt}
\end{center}

% ============================================================
\section*{Context: Pemex by the Numbers (2025--2026)}
% ============================================================

\begin{tcolorbox}[colback=white, colframe=black]
\small
\begin{center}
\begin{tabular}{ll}
\toprule
\textbf{Metric} & \textbf{Data} \\
\midrule
Financial debt (Dec 2025) & \$84.5 billion \\
Debt to suppliers/contractors & $\sim$\$20 billion more \\
Combined total liabilities & $\sim$\$100--105 billion \\
Annual loss (2025) & \$2.6 billion \\
Government bailout (2025 alone) & $>$\$40 billion \\
Oil production & $\sim$1.6--1.8 million bpd (down 50\% from peak) \\
Debt as \% of Mexico's GDP & $\sim$8\% \\
Status & World's most indebted oil company \\
\bottomrule
\end{tabular}
\end{center}
\textbf{Sources:} Bloomberg, US News, World Oil, OilPrice.com (Feb 2026)
\end{tcolorbox}

% ============================================================
\section*{Grammar: Relative Clauses --- Deep Dive}
% ============================================================

\subsection*{Defining vs.\ Non-Defining}

\begin{itemize}[nosep]
\item \textbf{Defining} (no commas) = essential info, identifies which one:\\
``The president \textbf{who} expropriated oil was C\'ardenas.''
(Which president? The one who expropriated.)
\item \textbf{Non-defining} (WITH commas) = extra info, could be removed:\\
``C\'ardenas, \textbf{who} served from 1934--1940, expropriated oil.''
(We already know which C\'ardenas.)
\end{itemize}

\textbf{TOEFL Rule:} ``That'' is ONLY used in defining clauses.
NEVER: ``C\'ardenas, \textbf{that} served...'' $\to$ WRONG.

\subsection*{Reduced Relative Clauses (very frequent in TOEFL)}

\begin{center}
\begin{tabular}{ll}
\toprule
\textbf{Full} & \textbf{Reduced} \\
\midrule
The company \textbf{that was created} in 1938 & The company \textbf{created} in 1938 \\
Workers \textbf{who are employed} by Pemex & Workers \textbf{employed} by Pemex \\
The debt \textbf{which is owed} to suppliers & The debt \textbf{owed} to suppliers \\
Pipelines \textbf{that are being tapped} & Pipelines \textbf{being tapped} \\
\bottomrule
\end{tabular}
\end{center}

\textbf{Rule:} Remove \texttt{who/which/that + be}. Keep the participle.

\subsection*{Whose, Where, When}

\begin{itemize}[nosep]
\item \textbf{Whose} = possession: ``Pemex, \textbf{whose} debt
exceeds \$84 billion, is the world's most indebted oil company.''
\item \textbf{Where} = place: ``The Balsas region, \textbf{where}
oil was first extracted, remains significant.''
\item \textbf{When} = time: ``1938, \textbf{when} C\'ardenas
expropriated oil, is celebrated as a national holiday.''
\end{itemize}

% ============================================================
\section*{Reading Passage 1: The 1938 Expropriation ($\sim$240 words)}
% ============================================================

\begin{tcolorbox}[colback=white, colframe=black]
\small
On March 18, 1938, President L\'azaro C\'ardenas announced the
expropriation of all foreign oil companies operating in Mexico. The
decision, which followed years of labor disputes between Mexican
workers and American and British petroleum firms, was grounded in
Article 27 of the 1917 Constitution, which asserted the nation's
sovereignty over its subsoil resources.

The reaction was extraordinary. Mexicans from every social class
contributed money, jewelry, and even livestock to help pay
compensation to the foreign companies. The expropriation became a
defining moment of national identity --- proof that a developing
nation could stand up to imperial economic powers.

Was it a mistake? The answer depends on the timeframe. In the short
and medium term, the expropriation was a triumph. Pemex became the
engine of Mexico's postwar development, funding infrastructure,
education, and social programs. Oil revenues constituted up to 40\%
of the federal budget at their peak.

In the long term, however, the state monopoly created conditions for
chronic problems: political interference in management decisions,
a bloated workforce protected by one of Mexico's most powerful
unions, systematic corruption at every level, and a dangerous
dependence on oil revenues that left the national budget vulnerable
to global price fluctuations.

C\'ardenas was not wrong to expropriate. But the generations that
followed were wrong to treat Pemex as an infinite piggy bank rather
than a company that required investment, modernization, and
accountability.
\end{tcolorbox}

\subsection*{Questions}

\begin{enumerate}[label=\textbf{\arabic*.}]
\item The main argument of the passage is:
\begin{enumerate}[label=(\Alph*)]
\item The expropriation was always wrong
\item The expropriation was right, but subsequent mismanagement
created the crisis
\item C\'ardenas should have given oil to foreign companies
\item Pemex was always profitable
\end{enumerate}

\item The word ``bloated'' most likely means:
\begin{enumerate}[label=(\Alph*)]
\item efficient \item excessively large and inefficient
\item small \item modern
\end{enumerate}

\item ``Piggy bank'' is used metaphorically to mean:
\begin{enumerate}[label=(\Alph*)]
\item A children's toy
\item A source of money used without concern for sustainability
\item A foreign investment fund
\item A bank for farmers
\end{enumerate}

\item Article 27 of the Constitution establishes:
\begin{enumerate}[label=(\Alph*)]
\item That oil belongs to foreign companies
\item National sovereignty over subsoil resources
\item That unions control all companies
\item Free trade agreements
\end{enumerate}

\item It can be inferred that the author:
\begin{enumerate}[label=(\Alph*)]
\item Blames C\'ardenas entirely for Pemex's problems
\item Distinguishes between the original decision and later mismanagement
\item Thinks Mexico should never have had oil
\item Supports unlimited government spending on Pemex
\end{enumerate}
\end{enumerate}

\textbf{Answers:} 1-B, 2-B, 3-B, 4-B, 5-B

% ============================================================
\section*{Reading Passage 2: Huachicol \& The Modern Crisis ($\sim$240 words)}
% ============================================================

\begin{tcolorbox}[colback=white, colframe=black]
\small
In recent years, the word ``huachicol'' --- originally slang for
adulterated alcohol --- has become synonymous with one of Mexico's
most destructive criminal enterprises: the systematic theft of fuel
from Pemex pipelines. Criminal organizations, often working with
complicit Pemex employees, tap into pipelines across the country,
extracting millions of barrels of fuel annually.

The scale is staggering. According to the Financial Times, fuel
theft costs Pemex an estimated \$3 billion per year. But the
economic damage is only part of the story. In January 2019, a
ruptured pipeline in Tlahuelilpan, Hidalgo, exploded while
villagers were collecting leaked gasoline, killing 137 people ---
including children. The tragedy illustrated how poverty, corruption,
and institutional failure converge in the huachicol crisis.

The problem has deeper roots than simple criminality. Decades of
underinvestment in pipeline security, combined with salaries so low
that corruption becomes a survival strategy for workers, created an
environment where fuel theft was almost inevitable. Some analysts
describe it as a symptom of what they call ``institutional entropy''
--- the gradual decay of a system that was never adequately
maintained or reformed.

Meanwhile, Pemex's total obligations exceed \$100 billion. Its
oil production has fallen nearly 50\% from peak levels. The
government continues to pour tens of billions into rescue packages,
effectively transferring Pemex's debt onto the shoulders of Mexican
taxpayers --- the same people C\'ardenas sought to empower in 1938.
The irony could not be more Pynchonesque.
\end{tcolorbox}

\subsection*{Questions}

\begin{enumerate}[label=\textbf{\arabic*.}, start=6]
\item ``Huachicol'' originally referred to:
\begin{enumerate}[label=(\Alph*)]
\item Stolen fuel \item Adulterated alcohol
\item A type of pipeline \item A criminal organization
\end{enumerate}

\item The Tlahuelilpan tragedy illustrates:
\begin{enumerate}[label=(\Alph*)]
\item That fuel theft is harmless
\item The deadly convergence of poverty, corruption, and institutional failure
\item Successful government intervention
\item Pipeline safety improvements
\end{enumerate}

\item ``Institutional entropy'' in this context means:
\begin{enumerate}[label=(\Alph*)]
\item Rapid growth \item Gradual systemic decay from neglect
\item A physics concept only \item Government reform
\end{enumerate}

\item The final sentence's ``irony'' refers to:
\begin{enumerate}[label=(\Alph*)]
\item C\'ardenas would have approved of the current situation
\item The people C\'ardenas sought to help now bear Pemex's debt burden
\item Pemex is profitable
\item Pynchon wrote about Pemex
\end{enumerate}

\item The word ``complicit'' is closest in meaning to:
\begin{enumerate}[label=(\Alph*)]
\item innocent \item involved in wrongdoing / knowingly participating
\item unaware \item heroic
\end{enumerate}
\end{enumerate}

\textbf{Answers:} 6-B, 7-B, 8-B, 9-B, 10-B

% ============================================================
\section*{Relative Clause Exercises (10 questions)}
% ============================================================

\begin{enumerate}[label=\textbf{\arabic*.}, leftmargin=*]

\item L\'azaro C\'ardenas, \rule{2cm}{0.4pt} expropriated oil in
1938, is considered a national hero.
\begin{enumerate}[label=(\Alph*)]
\item which \item that \item who \item whose
\end{enumerate}

\item The pipelines \rule{2cm}{0.4pt} by cartels are spread across
the entire country.
\begin{enumerate}[label=(\Alph*)]
\item tapping \item tapped \item which tapping \item are tap
\end{enumerate}

\item Pemex, \rule{2cm}{0.4pt} debt exceeds \$84 billion, is the
world's most indebted oil company.
\begin{enumerate}[label=(\Alph*)]
\item which \item who \item whose \item that
\end{enumerate}

\item The explosion \rule{2cm}{0.4pt} killed 137 people occurred in
Tlahuelilpan.
\begin{enumerate}[label=(\Alph*)]
\item who \item that \item whose \item where
\end{enumerate}

\item March 18, 1938, \rule{2cm}{0.4pt} the expropriation was
announced, is celebrated annually in Mexico.
\begin{enumerate}[label=(\Alph*)]
\item which \item that \item when \item where
\end{enumerate}

\item \textit{Find error:}
Pemex, (A)\textbf{that} was created in 1938, (B)\textbf{has become}
the (C)\textbf{most indebted} oil company (D)\textbf{in} the world.

\item Workers \rule{2cm}{0.4pt} salaries are too low often turn to
corruption as a survival strategy.
\begin{enumerate}[label=(\Alph*)]
\item who \item whose \item which \item whom
\end{enumerate}

\item The region \rule{2cm}{0.4pt} fuel theft is most common lacks
basic infrastructure.
\begin{enumerate}[label=(\Alph*)]
\item which \item who \item where \item whose
\end{enumerate}

\item \textit{Find error:}
The people (A)\textbf{whom} C\'ardenas sought to (B)\textbf{empower}
now (C)\textbf{bears} the burden of (D)\textbf{Pemex's} debt.

\item The government's rescue packages, \rule{2cm}{0.4pt} total
billions of dollars, have failed to solve the structural crisis.
\begin{enumerate}[label=(\Alph*)]
\item that \item who \item which \item whose
\end{enumerate}
\end{enumerate}

\subsection*{Answer Key}

\begin{center}
\begin{tabular}{cll}
\toprule
\# & Ans & Explanation \\
\midrule
1 & C & ``Who'' for people (C\'ardenas). Non-defining $\to$ no ``that'' \\
2 & B & Reduced passive: ``(that are) tapped by cartels'' \\
3 & C & ``Whose'' = possessive (Pemex's debt) \\
4 & B & ``That'' for defining clause (identifies which explosion) \\
5 & C & ``When'' for time references \\
6 & A & ``That'' $\to$ ``which'' (non-defining with commas) \\
7 & B & ``Whose'' = possession (their salaries) \\
8 & C & ``Where'' for places \\
9 & C & ``People...BEARS'' $\to$ ``BEAR'' (people = plural) \\
10 & C & ``Which'' for non-defining clause about packages \\
\bottomrule
\end{tabular}
\end{center}

% ============================================================
\section*{Vocabulary from Today's Readings}
% ============================================================

\begin{center}
\small
\begin{tabular}{lll}
\toprule
\textbf{Word} & \textbf{Meaning (ES)} & \textbf{TOEFL usage} \\
\midrule
expropriation & expropiaci\'on & government seizure of property \\
sovereignty & soberan\'ia & national sovereignty \\
subsoil & subsuelo & subsoil resources \\
complicit & c\'omplice & complicit in wrongdoing \\
convergence & convergencia & the convergence of factors \\
chronic & cr\'onico & chronic problems \\
bloated & inflado/hinchado & a bloated bureaucracy \\
adulterated & adulterado & adulterated products \\
staggering & asombroso & staggering amounts \\
entropy & entrop\'ia & institutional entropy \\
\bottomrule
\end{tabular}
\end{center}

% ============================================================
\section*{Shadowing (Gemini Live --- 15 min)}
% ============================================================

\begin{quote}\small
``Criminal organizations, often working with complicit Pemex
employees, tap into pipelines across the country, extracting
millions of barrels of fuel annually. The scale is staggering.
Decades of underinvestment in pipeline security, combined with
salaries so low that corruption becomes a survival strategy,
created an environment where fuel theft was almost inevitable.''
\end{quote}

\textbf{Targets:} \textbf{expropriation} (ex-PRO-pree-AY-shun),
\textbf{sovereignty} (SOV-rin-tee), \textbf{complicit}
(com-PLIS-it), \textbf{staggering} (STAG-er-ing),
\textbf{inevitable} (in-EV-i-ta-bul).

% ============================================================
\section*{Speaking Bonus (Gemini Live)}
% ============================================================

\begin{tcolorbox}[colback=white, colframe=black]
Tell Gemini Live:

\textit{``I want to practice arguing a position in English. My topic
is: Was the 1938 oil expropriation in Mexico a mistake? I will
argue both sides for 2 minutes each. Correct my grammar and suggest
better vocabulary after each argument.''}

\textbf{Useful structures:}
\begin{itemize}[nosep]
\item ``On one hand... On the other hand...''
\item ``While it is true that..., it must also be noted that...''
\item ``The evidence suggests that...''
\item ``In retrospect, the decision was justified because...''
\item ``However, the long-term consequences indicate that...''
\end{itemize}
\end{tcolorbox}

\end{document}